\documentclass[11pt]{article}

\usepackage{amsmath, amssymb, amsfonts}
\usepackage{physics}
\usepackage{graphicx}
\usepackage{geometry}
\usepackage{siunitx}
\usepackage{bm}
\geometry{margin=1in}

\title{A Mathematical Tutorial on Ballistics for Precision Rifle Shooting}
\author{}
\date{}

\begin{document}
	
	\maketitle
	
	\tableofcontents
	
	\newpage
	
	\section{Introduction}
	
	Precision rifle shooting relies on understanding the physics governing projectile motion. Unlike idealized textbook problems, real-world ballistics involves:
	\begin{itemize}
		\item Gravity
		\item Aerodynamic drag
		\item Wind effects
		\item Bullet stability and spin
	\end{itemize}
	
	This tutorial develops the mathematical framework necessary to model bullet trajectories accurately.
	
	\section{Coordinate System and Notation}
	
	We define:
	\begin{itemize}
		\item $x$: horizontal distance (downrange)
		\item $y$: vertical height
		\item $z$: lateral deviation (wind direction)
	\end{itemize}
	
	Velocity vector:
	\[
	\bm{v} = (v_x, v_y, v_z)
	\]
	
	Acceleration:
	\[
	\bm{a} = \frac{d\bm{v}}{dt}
	\]
	
	\section{Basic Projectile Motion (Vacuum)}
	
	Ignoring air resistance:
	
	\[
	\bm{a} = (0, -g, 0)
	\]
	
	Solution:
	\begin{align}
		x(t) &= v_0 \cos\theta \cdot t \\
		y(t) &= v_0 \sin\theta \cdot t - \frac{1}{2}gt^2
	\end{align}
	
	Eliminating $t$:
	
	\[
	y(x) = x \tan\theta - \frac{g x^2}{2 v_0^2 \cos^2\theta}
	\]
	
	\subsection{Limitations}
	
	This model is insufficient because:
	\begin{itemize}
		\item Bullets travel at supersonic speeds
		\item Drag dominates trajectory behavior
	\end{itemize}
	
	\section{Aerodynamic Drag}
	
	\subsection{Drag Force Model}
	
	Drag force is:
	
	\[
	\bm{F}_D = -\frac{1}{2} \rho C_D A v \bm{v}
	\]
	
	where:
	\begin{itemize}
		\item $\rho$: air density
		\item $C_D$: drag coefficient
		\item $A$: cross-sectional area
		\item $v = \|\bm{v}\|$
	\end{itemize}
	
	Equation of motion:
	
	\[
	m \frac{d\bm{v}}{dt} = -\frac{1}{2} \rho C_D A v \bm{v} + (0, -mg, 0)
	\]
	
	\subsection{Ballistic Coefficient}
	
	Define ballistic coefficient:
	
	\[
	BC = \frac{m}{C_D A}
	\]
	
	Rewriting:
	
	\[
	\frac{d\bm{v}}{dt} = -\frac{\rho}{2 BC} v \bm{v} + (0, -g, 0)
	\]
	
	\section{1D Drag Model}
	
	Consider horizontal motion only:
	
	\[
	\frac{dv}{dt} = -k v^2, \quad k = \frac{\rho}{2 BC}
	\]
	
	Separable equation:
	
	\[
	\int \frac{dv}{v^2} = -k \int dt
	\]
	
	Solution:
	
	\[
	-\frac{1}{v} = -kt + C
	\]
	
	\[
	v(t) = \frac{v_0}{1 + k v_0 t}
	\]
	
	Position:
	
	\[
	x(t) = \frac{1}{k} \ln(1 + k v_0 t)
	\]
	
	\section{Full 2D System with Drag}
	
	We solve:
	
	\begin{align}
		\frac{dv_x}{dt} &= -k v v_x \\
		\frac{dv_y}{dt} &= -k v v_y - g
	\end{align}
	
	with:
	
	\[
	v = \sqrt{v_x^2 + v_y^2}
	\]
	
	This system has no closed-form solution and must be solved numerically.
	
	\section{Numerical Integration}
	
	\subsection{Euler Method}
	
	\begin{align}
		v_x^{n+1} &= v_x^n + \Delta t (-k v^n v_x^n) \\
		v_y^{n+1} &= v_y^n + \Delta t (-k v^n v_y^n - g)
	\end{align}
	
	\subsection{Runge-Kutta (RK4)}
	
	More accurate:
	
	\[
	\bm{v}_{n+1} = \bm{v}_n + \frac{\Delta t}{6}(k_1 + 2k_2 + 2k_3 + k_4)
	\]
	
	Used in modern ballistic solvers.
	
	\section{Wind Effects}
	
	Assume wind velocity $\bm{w}$.
	
	Relative velocity:
	
	\[
	\bm{v}_{rel} = \bm{v} - \bm{w}
	\]
	
	Drag becomes:
	
	\[
	\bm{F}_D = -\frac{1}{2} \rho C_D A \|\bm{v}_{rel}\| \bm{v}_{rel}
	\]
	
	\subsection{Crosswind Drift}
	
	Approximate lateral drift:
	
	\[
	z \approx \frac{w}{v_{avg}} x
	\]
	
	More accurate models require full 3D integration.
	
	\section{Spin Drift and Coriolis Effect}
	
	\subsection{Spin Drift}
	
	Due to gyroscopic stability:
	
	\[
	z_{spin} \propto x^2
	\]
	
	Typically small but significant at long range.
	
	\subsection{Coriolis Effect}
	
	Earth rotation introduces acceleration:
	
	\[
	\bm{a}_c = -2 \bm{\Omega} \times \bm{v}
	\]
	
	Magnitude:
	
	\[
	a_c \approx 2 \Omega v
	\]
	
	where $\Omega \approx 7.29 \times 10^{-5}$ rad/s.
	
	\section{Trajectory Correction}
	
	\subsection{Drop Compensation}
	
	Vertical drop:
	
	\[
	\Delta y \approx \frac{1}{2} g t^2
	\]
	
	Converted to angular correction:
	
	\[
	\theta \approx \frac{\Delta y}{x}
	\]
	
	\subsection{Minute of Angle (MOA)}
	
	\[
	1 \text{ MOA} \approx 1.047 \text{ inches at 100 yards}
	\]
	
	Correction:
	
	\[
	\text{MOA} = \frac{\Delta y}{x} \cdot \frac{180 \times 60}{\pi}
	\]
	
	\subsection{Milliradians (MIL)}
	
	\[
	1 \text{ mil} = 10^{-3} \text{ rad}
	\]
	
	\[
	\text{mil} = \frac{\Delta y}{x} \times 1000
	\]
	
	\section{Atmospheric Effects}
	
	\subsection{Air Density}
	
	\[
	\rho = \frac{P}{R T}
	\]
	
	where:
	\begin{itemize}
		\item $P$: pressure
		\item $T$: temperature
		\item $R$: gas constant
	\end{itemize}
	
	Higher altitude $\Rightarrow$ lower $\rho$ $\Rightarrow$ less drag.
	
	\section{Practical Ballistic Modeling}
	
	\subsection{G1 and G7 Models}
	
	Drag coefficient depends on velocity:
	
	\[
	C_D = f(v)
	\]
	
	Standard drag models:
	\begin{itemize}
		\item G1: flat-base bullets
		\item G7: boat-tail bullets
	\end{itemize}
	
	\subsection{Table-Based Integration}
	
	Modern solvers use:
	\begin{itemize}
		\item Tabulated drag curves
		\item Interpolation
		\item Numerical integration
	\end{itemize}
	
	\section{Example Calculation}
	
	Given:
	\begin{itemize}
		\item $v_0 = 800$ m/s
		\item $BC = 0.5$
		\item Distance = 500 m
	\end{itemize}
	
	Estimate time of flight (approximate):
	
	\[
	t \approx \frac{x}{v_{avg}} \approx \frac{500}{700} \approx 0.71 \text{ s}
	\]
	
	Drop:
	
	\[
	\Delta y = \frac{1}{2} g t^2 \approx 2.5 \text{ m}
	\]
	
	MIL correction:
	
	\[
	\text{mil} = \frac{2.5}{500} \times 1000 = 5 \text{ mil}
	\]
	
	\section{Conclusion}
	
	Precision shooting requires combining:
	\begin{itemize}
		\item Physics (drag, gravity)
		\item Mathematics (ODEs, numerical methods)
		\item Environmental corrections
	\end{itemize}
	
	Modern ballistic solvers implement these equations numerically for high accuracy.
	
\end{document}